% ══════════════════════════════════════════════════════════════════
%  FICHE D'EXERCICES — Chapitre 13 : Oxydoréduction — piles et électrolyse
% ══════════════════════════════════════════════════════════════════
\section{Exercices type Baccalauréat}
\textit{Données : constante de Faraday $F = \SI{96500}{\coulomb\per\mol}$ ;
$M(\ce{Cu})=\SI{63.5}{\gram\per\mol}$, $M(\ce{Zn})=\SI{65.4}{\gram\per\mol}$,
$M(\ce{Ag})=\SI{108}{\gram\per\mol}$.}

\rubrique{Couples redox, demi-équations}

\exo{}\diff{1} Pour chacun des couples, écrire la demi-équation électronique :
a)~$\ce{Cu^2+}/\ce{Cu}$ ; b)~$\ce{Fe^3+}/\ce{Fe^2+}$ ; c)~$\ce{MnO4-}/\ce{Mn^2+}$
(milieu acide) ; d)~$\ce{Cr2O7^2-}/\ce{Cr^3+}$ (milieu acide).

\exo{}\diff{2} Équilibrer l'équation de la réaction entre les ions permanganate et
les ions fer(II) en milieu acide :
$\ce{MnO4- + Fe^2+ + H+ -> Mn^2+ + Fe^3+ + H2O}$.

\exo{}\diff{2} Une lame de zinc plongée dans une solution de sulfate de cuivre se
recouvre de cuivre et la solution se décolore.
\begin{enumerate}[label=\arabic*)]
\item Identifier l'oxydant et le réducteur, écrire les deux demi-équations.
\item Écrire l'équation bilan de la réaction.
\end{enumerate}

\rubrique{Piles électrochimiques}

\exo{}\diff{2} On réalise la pile Daniell : électrode de zinc dans $\ce{Zn^2+}$,
électrode de cuivre dans $\ce{Cu^2+}$, reliées par un pont salin. La pile débite et
la masse de l'électrode de zinc diminue.
\begin{enumerate}[label=\arabic*)]
\item Faire un schéma annoté de la pile.
\item Identifier l'anode et la cathode, préciser leur polarité.
\item Écrire les demi-équations aux électrodes et l'équation bilan de fonctionnement.
\item Quel est le rôle du pont salin ?
\end{enumerate}

\exo{}\diff{2} Une pile fonctionne et débite un courant d'intensité constante
$I = \SI{0.20}{\ampere}$ pendant \SI{1}{\hour}30. À l'électrode de cuivre, il se
dépose du cuivre métallique ($\ce{Cu^2+ + 2 e- -> Cu}$).
\begin{enumerate}[label=\arabic*)]
\item Calculer la quantité d'électricité débitée.
\item En déduire la quantité de matière puis la masse de cuivre déposée.
\end{enumerate}

\rubrique{Électrolyse}

\exo{}\diff{1} Définir l'électrolyse. En quoi diffère-t-elle du fonctionnement d'une
pile (sens des transformations, énergie) ?

\exo{}\diff{2} On électrolyse une solution de chlorure de cuivre(II) avec des
électrodes inertes. Il se dépose du cuivre à la cathode et il se dégage du dichlore à
l'anode.
\begin{enumerate}[label=\arabic*)]
\item Écrire les demi-équations aux électrodes et l'équation bilan.
\item Calculer la masse de cuivre déposée si l'on fait passer un courant
$I = \SI{1.0}{\ampere}$ pendant \SI{30}{\minute}.
\end{enumerate}

\exo{}\diff{2} On veut argenter un objet par électrolyse d'une solution contenant des
ions $\ce{Ag+}$ ($\ce{Ag+ + e- -> Ag}$). On souhaite déposer \SI{1.08}{\gram}
d'argent.
\begin{enumerate}[label=\arabic*)]
\item Calculer la quantité de matière d'argent à déposer.
\item En déduire la quantité d'électricité nécessaire.
\item Quelle est la durée de l'électrolyse sous un courant $I = \SI{0.50}{\ampere}$ ?
\end{enumerate}

\rubrique{Problème de synthèse — type Baccalauréat}

\sujetbac[d'après Baccalauréat C, Cameroun]
On réalise une pile en associant les couples $\ce{Cu^2+}/\ce{Cu}$ et
$\ce{Ag+}/\ce{Ag}$. L'expérience montre que l'électrode de cuivre est le pôle négatif.
\begin{enumerate}[label=\arabic*)]
\item Écrire les demi-équations aux électrodes (en précisant oxydation/réduction) et
l'équation bilan de la réaction qui se produit lorsque la pile débite.
\item La pile débite un courant $I = \SI{50}{\milli\ampere}$ pendant \SI{2}{\hour}.
Calculer la quantité d'électricité débitée.
\item En déduire la masse de cuivre consommée et la masse d'argent déposée.
\item De combien varie la masse de chaque électrode ?
\end{enumerate}
